% ===== PRÉAMBULE CANONIQUE — à coller VERBATIM en tête de CHAQUE .tex =====
\documentclass[11pt,a4paper]{article}
\usepackage[utf8]{inputenc}\usepackage[T1]{fontenc}\usepackage{lmodern}
\usepackage[french]{babel}
\usepackage{amsmath,amssymb,mathtools}\usepackage{icomma}
\usepackage{siunitx}\sisetup{locale=FR}  % \qty{}{}, \si{}, \ang{}
\usepackage{xcolor}\usepackage{colortbl}
\usepackage{tikz}\usepackage{pgfplots}\pgfplotsset{compat=1.18}
\usepackage[siunitx,europeanresistors]{circuitikz} % circuits (élec.)
\usepackage[most]{tcolorbox}\usepackage{enumitem}\usepackage{array,booktabs}
\usepackage[a4paper,margin=2.2cm]{geometry}
\usetikzlibrary{arrows.meta,calc,angles,quotes,positioning,patterns,%
  backgrounds,shapes.geometric,decorations.pathmorphing,decorations.markings,intersections}
\definecolor{primary}{RGB}{222,49,99}\definecolor{primarydark}{RGB}{150,22,55}
\definecolor{defcol}{RGB}{0,114,178}\definecolor{methcol}{RGB}{52,92,148}
\definecolor{excol}{RGB}{0,135,90}\definecolor{attncol}{RGB}{200,40,55}
\tcbset{boxbase/.style={enhanced,breakable,boxrule=0.9pt,arc=2.5pt,
  left=3mm,right=3mm,top=2.3mm,bottom=2.3mm,fonttitle=\bfseries,coltitle=white}}
% --- boîtes TUTORIEL (noms EXACTS, ne pas renommer) ---
\newtcolorbox{cle}[1][]{boxbase,colback=defcol!6,colframe=defcol,
  title={\textsc{À retenir}\ifx\relax#1\relax\relax\else\textmd{ : \itshape #1}\fi}}
\newtcolorbox{methode}[1][]{boxbase,colback=methcol!7,colframe=methcol,sharp corners,
  title={\textsc{Méthode}\ifx\relax#1\relax\relax\else\textmd{ --- \itshape #1}\fi}}
\newtcolorbox{exemplebox}[1][]{boxbase,colback=excol!5,colframe=excol,
  title={\textsc{Exemple}\ifx\relax#1\relax\relax\else\textmd{ : \itshape #1}\fi}}
\newtcolorbox{piege}[1][]{boxbase,colback=attncol!6,colframe=attncol,
  title={\textsc{Piège}\ifx\relax#1\relax\relax\else\textmd{ : \itshape #1}\fi}}
\newtcolorbox{remarque}[1][]{boxbase,colback=black!4,colframe=black!45,
  title={\textsc{Remarque}\ifx\relax#1\relax\relax\else\textmd{ : \itshape #1}\fi}}
% --- boîtes EXERCICE / CORRIGÉ (noms EXACTS, auto-comptées) ---
\newtcolorbox[auto counter]{exo}[1][]{boxbase,colback=primary!4,colframe=primary,
  title={\textsc{Exercice}~\thetcbcounter\ifx\relax#1\relax\relax\else\textmd{ --- \itshape #1}\fi}}
\newtcolorbox[auto counter]{corrige}[1][]{boxbase,colback=excol!4,colframe=excol,
  title={\textsc{Corrigé}~\thetcbcounter\ifx\relax#1\relax\relax\else\textmd{ --- \itshape #1}\fi}}
\newcommand{\vv}[1]{\vec{#1}}\newcommand{\ds}{\displaystyle}
\newcommand{\R}{\mathbb{R}}\newcommand{\N}{\mathbb{N}}\newcommand{\Z}{\mathbb{Z}}
% (un \usepackage{fancyhdr}+en-tête est possible ; il est ignoré côté web)
% =========================================================================
\begin{document}
\sisetup{per-mode=symbol}

\begin{corrige}[retrouver la hauteur de chute]
\begin{enumerate}[label=\alph*)]
  \item $v=\sqrt{2gh}\Rightarrow h=\dfrac{v^2}{2g}=\dfrac{40^2}{2\times10}=\dfrac{1600}{20}=\qty{80}{\metre}$.
  \item $t=\dfrac{v}{g}=\dfrac{40}{10}=\qty{4}{\second}$.
\end{enumerate}
\end{corrige}

\begin{corrige}[lancer vertical vers le haut]
Axe vers le haut, $a=-g$.
\begin{enumerate}[label=\alph*)]
  \item Au sommet $v=0$ : $t_s=\dfrac{v_0}{g}=\dfrac{20}{10}=\qty{2}{\second}$.
  \item $y_{\max}=\dfrac{v_0^2}{2g}=\dfrac{400}{20}=\qty{20}{\metre}$.
  \item Par symétrie, la descente dure autant que la montée : durée totale $=2\,t_s=\qty{4}{\second}$.
\end{enumerate}
\end{corrige}

\begin{corrige}[tir horizontal complet]
\begin{enumerate}[label=\alph*)]
  \item Chute verticale : $t=\sqrt{\dfrac{2h}{g}}=\sqrt{\dfrac{2\times80}{10}}=\sqrt{16}=\qty{4}{\second}$.
  \item Portée $x=v_{0x}\,t=20\times4=\qty{80}{\metre}$.
  \item $v_x=\qty{20}{\metre\per\second}$, $v_y=g\,t=40\ \si{\metre\per\second}$ : $v=\sqrt{20^2+40^2}=\sqrt{2000}\approx\qty{44,7}{\metre\per\second}$.
\end{enumerate}
\end{corrige}

\begin{corrige}[la chute seconde après seconde]
\begin{enumerate}[label=\alph*)]
  \item $v=g\,t$ : \quad \qty{10}{}, \qty{20}{}, \qty{30}{}, \qty{40}{\metre\per\second}.
  \item $h=\tfrac12 g t^2$ : \quad \qty{5}{}, \qty{20}{}, \qty{45}{}, \qty{80}{\metre}.
  \item Entre $t=2$ et $t=\qty{3}{\second}$ : $45-20=\qty{25}{\metre}$ (les distances par seconde croissent en $5,15,25,35,\dots$).
\end{enumerate}
\end{corrige}

\begin{corrige}[hauteur maximale d'un lancer]
\begin{enumerate}[label=\alph*)]
  \item $y_{\max}=\dfrac{v_0^2}{2g}\Rightarrow v_0=\sqrt{2g\,y_{\max}}=\sqrt{2\times10\times45}=\sqrt{900}=\qty{30}{\metre\per\second}$.
  \item Torricelli (montée, $a=-g$) : $v^2=v_0^2-2g h=900-2\times10\times20=500 \Rightarrow v=\sqrt{500}\approx\qty{22,4}{\metre\per\second}$.
\end{enumerate}
\end{corrige}

\end{document}
